\section{Electric and Magnetic Fields from Potentials}
\label{sec:fields_from_potentials}

The definitions

\[
\mathbf{B}
=
\nabla\times\mathbf{A}
\]

and

\[
\mathbf{E}
=
-\nabla\Phi
-
\frac{\partial\mathbf{A}}{\partial t}
\]

are designed so that the homogeneous Maxwell equations hold identically. This section verifies those identities, examines their component form, and applies the formulas to representative configurations.

\subsection{Magnetic Gauss law}

Using

\[
\mathbf{B}
=
\nabla\times\mathbf{A},
\]

we obtain

\[
\nabla\cdot\mathbf{B}
=
\nabla\cdot
\left(
\nabla\times\mathbf{A}
\right).
\]

The divergence of a curl vanishes for every twice continuously differentiable vector field:

\[
\boxed{
\nabla\cdot\mathbf{B}=0
}.
\]

In Cartesian components,

\[
B_i
=
\epsilon_{ijk}
\partial_jA_k.
\]

Therefore

\[
\partial_iB_i
=
\epsilon_{ijk}
\partial_i\partial_jA_k.
\]

The derivative factor is symmetric under

\[
i\longleftrightarrow j,
\]

while \(\epsilon_{ijk}\) is antisymmetric. Their contraction vanishes:

\[
\epsilon_{ijk}
\partial_i\partial_jA_k
=0.
\]

The identity follows from the commutativity of ordinary partial derivatives.

\subsection{Faraday's law}

Take the curl of the electric field:

\[
\begin{aligned}
\nabla\times\mathbf{E}
&=
-\nabla\times\nabla\Phi
-
\nabla\times
\frac{\partial\mathbf{A}}{\partial t}.
\end{aligned}
\]

The curl of a gradient vanishes:

\[
\nabla\times\nabla\Phi=0.
\]

For sufficiently regular fields, time and spatial derivatives commute:

\[
\nabla\times
\frac{\partial\mathbf{A}}{\partial t}
=
\frac{\partial}{\partial t}
\left(
\nabla\times\mathbf{A}
\right)
=
\frac{\partial\mathbf{B}}{\partial t}.
\]

Thus

\[
\boxed{
\nabla\times\mathbf{E}
+
\frac{\partial\mathbf{B}}{\partial t}
=0
}.
\]

Faraday's law is therefore built into the potential representation.

\subsection{Component verification}

Write

\[
E_i
=
-\partial_i\Phi
-
\frac{\partial A_i}{\partial t}
\]

and

\[
B_i
=
\epsilon_{ijk}\partial_jA_k.
\]

The \(i\)-th component of the curl of \(\mathbf{E}\) is

\[
\begin{aligned}
(\nabla\times\mathbf{E})_i
&=
\epsilon_{ijk}\partial_jE_k\\
&=
-\epsilon_{ijk}
\partial_j\partial_k\Phi
-
\epsilon_{ijk}
\partial_j
\frac{\partial A_k}{\partial t}.
\end{aligned}
\]

The first term vanishes by symmetry and antisymmetry. The second term is

\[
-
\frac{\partial}{\partial t}
\left(
\epsilon_{ijk}\partial_jA_k
\right)
=
-
\frac{\partial B_i}{\partial t}.
\]

Hence

\[
(\nabla\times\mathbf{E})_i
=
-
\frac{\partial B_i}{\partial t}.
\]

\subsection{Which Maxwell equations remain dynamical?}

The potential definitions automatically satisfy

\[
\nabla\cdot\mathbf{B}=0
\]

and

\[
\nabla\times\mathbf{E}
+
\frac{\partial\mathbf{B}}{\partial t}=0.
\]

The remaining equations,

\[
\nabla\cdot\mathbf{E}
=
\frac{\rho}{\varepsilon_0}
\]

and

\[
\nabla\times\mathbf{B}
-
\frac{1}{c^2}
\frac{\partial\mathbf{E}}{\partial t}
=
\mu_0\mathbf{J},
\]

determine how the potentials respond to electric charge and current. They will be derived from an action in the next chapter.

\subsection{Electrostatic limit}

If the potentials are time independent,

\[
\frac{\partial\mathbf{A}}{\partial t}=0,
\]

then

\[
\mathbf{E}
=
-\nabla\Phi.
\]

The electric field is curl free:

\[
\nabla\times\mathbf{E}=0.
\]

Gauss's law becomes a Poisson equation for the scalar potential:

\[
\boxed{
\nabla^2\Phi
=
-
\frac{\rho}{\varepsilon_0}
}.
\]

In a source-free region,

\[
\nabla^2\Phi=0.
\]

\subsection{Magnetostatic limit}

For steady currents,

\[
\frac{\partial\mathbf{E}}{\partial t}=0.
\]

Ampere's law is

\[
\nabla\times\mathbf{B}
=
\mu_0\mathbf{J}.
\]

Substitute

\[
\mathbf{B}
=
\nabla\times\mathbf{A}.
\]

Using the vector identity

\[
\nabla\times
\left(
\nabla\times\mathbf{A}
\right)
=
\nabla
\left(
\nabla\cdot\mathbf{A}
\right)
-
\nabla^2\mathbf{A},
\]

we obtain

\[
\nabla
\left(
\nabla\cdot\mathbf{A}
\right)
-
\nabla^2\mathbf{A}
=
\mu_0\mathbf{J}.
\]

In the Coulomb gauge,

\[
\nabla\cdot\mathbf{A}=0,
\]

this reduces to

\[
\boxed{
\nabla^2\mathbf{A}
=
-\mu_0\mathbf{J}
}.
\]

\subsection{Example: time-dependent uniform vector potential}

Let

\[
\Phi=0
\]

and

\[
\mathbf{A}(t)
=
A_0(t)\mathbf{e}_x.
\]

Because \(\mathbf{A}\) is spatially uniform,

\[
\mathbf{B}
=
\nabla\times\mathbf{A}
=0.
\]

The electric field is

\[
\boxed{
\mathbf{E}
=
-
\frac{\mathrm{d}A_0}{\mathrm{d}t}
\mathbf{e}_x
}.
\]

A time-dependent vector potential can therefore produce an electric field even when the magnetic field vanishes.

\subsection{Example: a plane-wave potential}

Choose

\[
\Phi=0
\]

and

\[
\mathbf{A}(z,t)
=
A_0
\cos(kz-\omega t)
\mathbf{e}_x.
\]

The electric field is

\[
\begin{aligned}
\mathbf{E}
&=
-
\frac{\partial\mathbf{A}}{\partial t}\\
&=
-
A_0\omega
\sin(kz-\omega t)
\mathbf{e}_x.
\end{aligned}
\]

The magnetic field is

\[
\begin{aligned}
\mathbf{B}
&=
\nabla\times\mathbf{A}\\
&=
-A_0k
\sin(kz-\omega t)
\mathbf{e}_y.
\end{aligned}
\]

For an electromagnetic wave in vacuum,

\[
\omega=ck.
\]

Therefore the magnitudes satisfy

\[
\boxed{
|\mathbf{E}|=c|\mathbf{B}|
}.
\]

The vectors \(\mathbf{E}\), \(\mathbf{B}\), and the propagation direction \(\mathbf{e}_z\) are mutually perpendicular. Their signs depend on the chosen phase and orientation of the potential.

\subsection{Example: time-dependent magnetic field}

Consider the symmetric-gauge potential

\[
\mathbf{A}
=
\frac{1}{2}
B(t)
\left(
-y\mathbf{e}_x
+x\mathbf{e}_y
\right)
\]

with

\[
\Phi=0.
\]

Its magnetic field is

\[
\mathbf{B}
=
B(t)\mathbf{e}_z.
\]

The induced electric field is

\[
\boxed{
\mathbf{E}
=
\frac{1}{2}
\frac{\mathrm{d}B}{\mathrm{d}t}
\left(
y\mathbf{e}_x
-
x\mathbf{e}_y
\right)
}.
\]

In cylindrical coordinates,

\[
\mathbf{E}
=
-
\frac{r}{2}
\frac{\mathrm{d}B}{\mathrm{d}t}
\mathbf{e}_{\varphi}.
\]

The electric field circulates around the changing magnetic field, as required by Faraday's law.

\subsection{Integral form of Faraday's law}

Integrate the differential equation over a surface \(S\):

\[
\int_S
\left(
\nabla\times\mathbf{E}
\right)
\cdot\mathbf{n}
\,\mathrm{d}A
=
-
\frac{\mathrm{d}}{\mathrm{d}t}
\int_S
\mathbf{B}\cdot\mathbf{n}
\,\mathrm{d}A.
\]

By Stokes' theorem,

\[
\boxed{
\oint_{\partial S}
\mathbf{E}\cdot\mathrm{d}\boldsymbol{\ell}
=
-
\frac{\mathrm{d}\Phi_B}{\mathrm{d}t}
}.
\]

Here

\[
\Phi_B
=
\int_S
\mathbf{B}\cdot\mathbf{n}
\,\mathrm{d}A
\]

is the magnetic flux through \(S\).

\subsection{Regularity assumptions}

The derivations used the commutativity of partial derivatives:

\[
\partial_i\partial_j
=
\partial_j\partial_i
\]

and

\[
\frac{\partial}{\partial t}
\partial_i
=
\partial_i
\frac{\partial}{\partial t}.
\]

These identities hold for sufficiently smooth potentials. Singular potentials, excluded regions, or distributional sources require additional care. The local equations remain valid away from singularities.

\subsection{Common mistakes}

Typical errors include:

\begin{enumerate}
    \item omitting the term \(-\partial\mathbf{A}/\partial t\) from the electric field;
    \item using \(\nabla\times\nabla\Phi\neq0\);
    \item forgetting that time and spatial derivatives were assumed to commute;
    \item concluding that \(\mathbf{B}=0\) implies \(\mathbf{A}=0\);
    \item using electrostatic formulas in a time-dependent induction problem;
    \item reversing the sign in Faraday's law.
\end{enumerate}

\subsection{What has been established}

The potential representation implies identically

\[
\nabla\cdot\mathbf{B}=0
\]

and

\[
\nabla\times\mathbf{E}
+
\frac{\partial\mathbf{B}}{\partial t}=0.
\]

The two homogeneous Maxwell equations are therefore structural consequences of expressing the fields in terms of \(\Phi\) and \(\mathbf{A}\). The next section analyzes the full family of potential transformations that leave both physical fields unchanged.
